% Irodov 用户底本 PDF 页92--94，原书印刷页83--85，原题185--194。
% 第185题底本印Δm=0.70g却简答1.2J；第189题简答mJ应为MJ。
\thermophaseproblem{185}{83}
在温度 $t=100^\circ\mathrm C$ 时，
无质量活塞下的圆筒内装有饱和水蒸气。
缓慢推动活塞使
$\Delta m=0.70\,\mathrm g$ 的蒸汽冷凝。
假设水蒸气是理想气体，
忽略水的液态体积，求对气体所做的功。
\exsolution
温度为 $T=373.15\,\mathrm K$；
饱和冷凝过程等温且近似等压
$p_{\rm sat}=p_0$。
冷凝物的气相物质的量
$\Delta\nu=\Delta m/M$，
其原先占据的气相体积为
$\Delta V_{\rm vap}=\Delta\nu RT/p_0$。
忽略液态体积，活塞压缩体积
$|\Delta V|=\Delta\nu RT/p_0$，
所以外界对气体所做的正功
\[
\boxed{A_{\rm on}=p_0|\Delta V|
=\frac{\Delta m RT}{M}}.
\]
取水摩尔质量
$M=18\times10^{-3}\,\mathrm{kg/mol}$，
题面质量 $\Delta m=7.0\times10^{-4}\,\mathrm{kg}$，
得
\[
\boxed{A_{\rm on}\simeq1.21\times10^2\,\mathrm J=120\,\mathrm J}.
\]
\excheck \textbf{底本数值冲突：}
书后简答印同一公式但给“$1.2\,\mathrm J$”。
使用题面清楚印出的 $0.70\,\mathrm g$，
由气体定律得到约 $120\,\mathrm J$，
相差100倍。未擅自把冷凝质量改为 $0.007\,\mathrm g$。

\thermophaseproblem{186}{84}
体积 $V=6.0\,\mathrm L$ 的容器里装有水和
饱和水蒸气，压强 $40\,\mathrm{atm}$、
温度 $250^\circ\mathrm C$。
气相比体积
$v_v'=50\,\mathrm{L/kg}$，
整个气液系统质量
$m=5.0\,\mathrm{kg}$。
求气相的质量及体积。
\exsolution
设气相质量 $m_v$，液相质量
$m_l=m-m_v$，
液态水的比体积记为
$v_l'\simeq1.0\,\mathrm{L/kg}$。
总容积为
\[
V=m_vv_v'+(m-m_v)v_l'
=m v_l'+m_v(v_v'-v_l').
\]
解得
\[
\boxed{m_v=\frac{V-mv_l'}{v_v'-v_l'}}.
\]
代入以升、千克计的数值：
\[
m_v=\frac{6.0-5.0(1.0)}{50-1.0}
\,\mathrm{kg}
\simeq\boxed{0.020\,\mathrm{kg}=20\,\mathrm g}.
\]
气相体积
\[
\boxed{V_v=m_vv_v'\simeq(0.0204)(50)
\,\mathrm L\simeq1.0\,\mathrm L}.
\]
\excheck 不能把全部 $6\,\mathrm L$
都当作气相容积；液相虽比体积小，
但总液体质量约5 kg，仍占约5 L。

\thermophaseproblem{187}{84}
温度 $100^\circ\mathrm C$ 时，
圆筒活塞下的饱和水蒸气初体积
$V_0=5.0\,\mathrm L$。
在等温条件下压缩至
$V=1.6\,\mathrm L$，
假设饱和蒸汽为理想气体，
求凝结成液体的质量。
\exsolution
饱和蒸汽在相变过程维持同温饱和压
$p_0\simeq1\,\mathrm{atm}$。
压缩减少的气相体积为
$\Delta V=V_0-V$；
忽略液相体积，凝结的摩尔数满足
$p_0\Delta V=\Delta\nu RT$。
因此
\[
\boxed{m_{\rm cond}=
\frac{Mp_0(V_0-V)}{RT}}.
\]
用 $M=0.018\,\mathrm{kg/mol}$、
$\Delta V=3.4\times10^{-3}\,\mathrm{m^3}$、
$T=373.15\,\mathrm K$：
\[
\boxed{m_{\rm cond}\simeq2.0\,\mathrm g}.
\]
\excheck 等温饱和压缩不是一团不凝结气体的
$pV=\text{常数}$ 过程；气相摩尔数随冷凝而减少。

\thermophaseproblem{188}{84}
一批饱和蒸汽最初只处于气相，
现等温将其所占总容积压缩 $n$ 倍。
假设饱和蒸汽比体积是液相比体积的
$N$ 倍，其中 $N>n$。
求末态液相占总容积的比例 $\eta$。
另若末态体积恰在
$p$--$V$ 等温相变平台的中点，求 $\eta$。
\exsolution
设液相比体积为 $v_l$，
气相比体积为 $v_v=Nv_l$，
物质总质量为 $m$。
最初全为饱和气相，
$V_0=mv_v$，
压缩后 $V_f=V_0/n$。
设末态液相质量 $m_l$，
气相质量 $m-m_l$：
\[
V_f=m_lv_l+(m-m_l)v_v
=mv_v-m_l(v_v-v_l).
\]
因此
$m_l=(V_0-V_f)/(v_v-v_l)$。
液相\emph{体积}占比
\[
\eta=\frac{m_lv_l}{V_f}
=\frac{(V_0-V_f)v_l}{V_f(v_v-v_l)}
=\boxed{\frac{n-1}{N-1}}.
\]
若末态体积在相变平台两端
$mv_l$ 与 $mv_v$ 的中点，
$V_f=m(v_l+v_v)/2$，
于是
\[
\boxed{\eta=\frac{mv_l}{2V_f}
=\frac1{N+1}}.
\]
\excheck 本题问的是\emph{液相体积占比}，
不是液态水的质量占比；
两者相差可很大。

\thermophaseproblem{189}{84}
质量 $m=1.00\,\mathrm{kg}$ 的水，
在标准大气压下已沸腾，
随后全部变为饱和水蒸气。
视饱和蒸汽为理想气体，
求系统的熵增与内能增量。
\exsolution
取沸点
$T=373.15\,\mathrm K$、
水的汽化潜热
$q\simeq2.26\times10^6\,\mathrm{J/kg}$。
等温定压相变吸热为
$Q=mq$，因此可逆相变的熵增
\[
\boxed{\Delta S=\frac{mq}{T}
\simeq\frac{2.26\times10^6}{373.15}
\simeq6.1\times10^3\,\mathrm{J/K}
=6.1\,\mathrm{kJ/K}}.
\]
由第一定律，
$\Delta U=Q-A$。
液相体积可忽略，
蒸汽的气体功
$A=p\Delta V\simeq\nu RT=mRT/M$。
于是
\[
\boxed{\Delta U=m\left(q-\frac{RT}{M}\right)}.
\]
以 $M=0.018\,\mathrm{kg/mol}$，
得
\[
\boxed{\Delta U\simeq2.1\times10^6\,\mathrm J
=2.1\,\mathrm{MJ}}.
\]
\excheck \textbf{底本单位勘误：}
本版书后印成“$2.1\,\mathrm{mJ}$”；
按题面1 kg水的汽化潜热和印出的内能公式，
必须是 $2.1\,\mathrm{MJ}$，与小写m的毫焦差 $10^9$ 倍。

\thermophaseproblem{190}{84}
质量 $m=20\,\mathrm g$、
初温 $0^\circ\mathrm C$ 的水，
装在绝热圆筒中，活塞无质量，
活塞面积
$S=410\,\mathrm{cm^2}$，
外界压强为标准大气压。
向水传入热量
$Q=20.0\,\mathrm{kJ}$，
求活塞上升高度。
\exsolution
活塞维持近似等压
$p_0=1\,\mathrm{atm}$。
先将水从 $0^\circ\mathrm C$
加热至 $100^\circ\mathrm C$，
需热
$Q_1=mc\Delta T$，
其中 $c\simeq4.18\times10^3\,\mathrm{J/(kg\,K)}$、
$\Delta T=100\,\mathrm K$。
\[
Q_1=(0.020)(4180)(100)\simeq8.36\,\mathrm{kJ}.
\]
余热
$Q_v=Q-Q_1\simeq11.64\,\mathrm{kJ}$
用于部分汽化。
取汽化潜热
$q\simeq2.26\,\mathrm{MJ/kg}$，
气相质量
$m_v=Q_v/q$。
在沸点 $T_b\simeq373\,\mathrm K$，
蒸汽近似理想气体且液相体积很小，
活塞体积增加
$\Delta V\simeq m_vRT_b/(Mp_0)$。
故
\[
\boxed{h\simeq
\frac{(Q-mc\Delta T)RT_b}{qMp_0S}}
\simeq\boxed{0.21\,\mathrm m\approx0.20\,\mathrm m}.
\]
\excheck 原书书后给约 $20\,\mathrm{cm}$；
采用常用表中汽化潜热为定压焓变时，
得到约21 cm。若把 $q$ 定义为内能汽化差，
还须计入气体膨胀功，
本稿明确了所用潜热约定。

\thermophaseproblem{191}{84}
无质量活塞下的绝热圆筒里有
$m_v=1.0\,\mathrm g$ 饱和水蒸气，
外界标准大气压。
往其中注入温度
$t_0=22^\circ\mathrm C$ 的
$m=1.0\,\mathrm g$ 冷水。
忽略活塞摩擦和筒热容，
求活塞下降时大气压所做的功。
\exsolution
初态蒸汽为 $100^\circ\mathrm C$，
系统仍在外压约1 atm下变化。
需要把新注入的水从
$T_0\simeq295\,\mathrm K$
加热到沸点
$T_b\simeq373\,\mathrm K$，
需热量
$Q_{\rm warm}=mc(T_b-T_0)$。
这部分热由一部分蒸汽冷凝释放，
故冷凝质量
\[
m_{\rm cond}=\frac{mc(T_b-T_0)}q.
\]
忽略液态体积，蒸汽体积减少
$\Delta V\simeq
(m_{\rm cond}/M)RT_b/p_0$。
大气压力在活塞下降过程中做正功
\[
\boxed{A=p_0\Delta V=
\frac{mc(T_b-T_0)RT_b}{qM}}.
\]
代入
$m=0.001\,\mathrm{kg}$、
$c=4180\,\mathrm{J/(kg\,K)}$、
$q\simeq2.26\times10^6\,\mathrm{J/kg}$，
\[
\boxed{A\simeq25\,\mathrm J}.
\]
\excheck 需要把冷水加热到与饱和蒸汽相同温度；
不能误认为整个原有1g蒸汽都冷凝。

\thermophaseproblem{192}{84}
已知球形凸液面上的饱和蒸汽压
比平面液面的饱和蒸汽压高
$\Delta p\ll p$，
并满足 Kelvin 近似公式
$\Delta p\simeq(\rho_v/\rho_l)(2\alpha/r)$。
在温度 $t=27^\circ\mathrm C$ 时，
求水滴直径 $d$，使
$\Delta p/p=\eta=1.0\%$；
假设蒸汽为理想气体。
\exsolution
用饱和蒸汽理想气体关系
$p=\rho_vRT/M$，
将 $\rho_v/p=M/(RT)$ 代入：
\[
\eta=\frac{\Delta p}{p}
\simeq\frac{\rho_v}{\rho_l}
\frac{2\alpha}{rp}
=\frac{2\alpha M}{r\rho_lRT}.
\]
水滴直径 $d=2r$，故
\[
\boxed{d\simeq
\frac{4\alpha M}{\eta\rho_lRT}}.
\]
取 $T=300.15\,\mathrm K$、
$\alpha\simeq0.072\,\mathrm{N/m}$、
$M=0.018\,\mathrm{kg/mol}$、
$\rho_l=1000\,\mathrm{kg/m^3}$，
\[
\boxed{d\simeq2.1\times10^{-7}\,\mathrm m
\simeq0.2\,\mathrm{\mu m}}.
\]
\excheck 公式只在蒸汽压相对增量很小的范围内适用；
更小水滴通常需使用非线性的 Kelvin 方程。

\thermophaseproblem{193}{84}
在 $t=100^\circ\mathrm C$ 的饱和水蒸汽下，
求每秒从 $1\,\mathrm{cm^2}$ 水面
离开的蒸汽分子总质量。
假设撞向水面的气相分子中，
仅有 $\eta=3.6\%$ 被液相留住。
\exsolution
饱和状态净蒸发质量通量为零，
因此离开液面的质量通量等于
气相分子被液体\emph{实际吸收}的通量。
饱和汽的撞击壁面质量通量为
\[
j_{\rm in}=\frac14\rho_v\bar v
=p\sqrt{\frac{M}{2\pi RT}},
\]
这是由 Maxwell 分布对法向速度积分得到的。
只保留其中比例 $\eta$，
净被吸收的质量通量为 $\eta j_{\rm in}$，
在动态平衡下同样等于离开液面的质量通量：
\[
\boxed{j_{\rm out}=
\eta p_0\sqrt{\frac{M}{2\pi RT}}}.
\]
用
$p_0=101325\,\mathrm{Pa}$、
$T=373.15\,\mathrm K$、
$M=0.018\,\mathrm{kg/mol}$，
计算
$j_{\rm out}\simeq3.5\,\mathrm{kg/(m^2\,s)}$。
换算 $1\,\mathrm{cm^2}=10^{-4}\,\mathrm{m^2}$：
\[
\boxed{\dot m\simeq3.5\times10^{-4}\,\mathrm{kg/s}
=0.35\,\mathrm{g/s}}
\quad\text{（每平方厘米水面）}.
\]
\excheck 这是饱和平衡下\emph{双向交换}
中的单向蒸发通量，不代表有同样大小的
净蒸发质量损失。

\thermophaseproblem{194}{85}
在温度 $T=2000\,\mathrm K$ 下，
钨丝在高真空中蒸发，
测得单位面积质量损失率
$\mu=1.2\times10^{-13}\,
\mathrm{g/(s\,cm^2)}$。
求该温度下钨的饱和蒸汽压强。
\exsolution
真空中的分子逸出通量
与饱和蒸汽压强之间的 Hertz--Knudsen
理想动力学近似为
\[
\mu=p_{\rm sat}\sqrt{\frac{M}{2\pi RT}}.
\]
故
\[
\boxed{p_{\rm sat}=
\mu\sqrt{\frac{2\pi RT}{M}}}.
\]
换算质量通量：
\[
1\,\mathrm{g/(cm^2\,s)}
=10\,\mathrm{kg/(m^2\,s)},
\quad
\mu=1.2\times10^{-12}\,\mathrm{kg/(m^2\,s)}.
\]
钨摩尔质量
$M\simeq0.184\,\mathrm{kg/mol}$，
因此
\[
\boxed{p_{\rm sat}\simeq
1.2\times10^{-12}
\sqrt{\frac{2\pi(8.314)(2000)}{0.184}}
\simeq9.0\times10^{-10}\,\mathrm{Pa}
=0.90\,\mathrm{nPa}}.
\]
\excheck 采用理想蒸发系数近似为1；
实际蒸发表面的凝结系数可能带来修正。
